\section{A Two-Engine Decentralized Data Platform}
\label{sec:twoengine}

The blockchain-native trajectory of Section~\ref{sec:chain} addresses analytics.
But analytics is only half of a data platform: the other half is the
transactional, per-user, per-project record store. These two workloads have
opposite shapes --- columnar scan-and-aggregate versus row-oriented
point-read-and-write --- and, crucially, opposite \emph{trust} requirements.
Aggregate analytics can run server-side (fast, shared, TEE-attested); a user's
own per-project records should be private end-to-end, visible to no operator at
all. We therefore propose two complementary engines over one decentralized
storage substrate, both native to the Lux network and both speaking ZAP.

\subsection{Engine 1: Datastore --- scalable analytics (OLAP)}

Datastore is the columnar engine of this paper: horizontally-scalable,
\emph{server-side} compute over the shared S-Chain object store
(Section~\ref{sec:chain}), demonstrated to scale out to multiple stateless
compute replicas sharing a single zero-copy copy of data
(Section~\ref{sec:eval}). Its privacy model is \emph{attestation}: compute runs
in a hardware trusted execution environment, so a client verifies an enclave
attestation rather than trusting the operator. This is the right model for
operational and aggregate analytics, where throughput and elastic scale dominate
and the data is the tenant's own aggregate.

\subsection{Engine 2: Decentralized SQL --- private per-user records (OLTP)}

The second engine is a SQLite-based transactional store --- one database file
per user or per project, the file being the unit of tenancy --- already running
under leaderless Quasar coordination in the sibling base-ha system. To this
foundation we add three primitives, each backed by an existing Lux chain, to
make it a \emph{decentralized, end-to-end-private} SQL engine:

\begin{description}[leftmargin=1.2em,style=nextline]
  \item[Client-held keys (K-Chain KEM).] Each user/project is issued a
  post-quantum key-encapsulation keypair from the K-Chain. The SQLite database
  --- at page or WAL granularity --- is encrypted under a data key wrapped by
  the user's KEM public key. The operator stores only ciphertext on S-Chain;
  \emph{decryption happens only on the client}. No operator, validator, or
  storage provider ever observes plaintext.
  \item[Compute on ciphertext (Z-Chain FHE).] To answer queries without
  decrypting, predicates and aggregates evaluate homomorphically against the
  encrypted pages using the Z-Chain's FHE/ZK primitives. The server returns an
  encrypted result the client alone can open. Where full FHE is too costly, the
  same chain's ZK path proves a result was computed correctly over committed
  ciphertext.
  \item[Coordinator-free merge (CRDTs).] Multi-device and multi-writer updates
  reconcile through conflict-free replicated data types, so concurrent writes
  merge deterministically without a consensus round on the write path. Quasar is
  used only for the coarse coordination (key rotation, file-version commitments
  on S-Chain), not for every row write.
\end{description}

The result is ``decentralized SQL'': a per-user record store that is private by
construction (E2E encryption), available without a trusted server (content on
S-Chain, keys on the client), and writable offline (CRDT merge), with the chain
providing only what chains are good at --- ordered commitments, key custody, and
verifiable compute.

\subsection{Why both, and how they compose}

The two engines are not redundant; they are the OLAP/OLTP split re-cast along the
trust axis (Table~\ref{tab:twoengine}). Datastore answers ``what are the
aggregate trends across the fleet'' at scale; Decentralized SQL holds ``this
user's project data'' privately. They share one substrate and one transport:

\begin{table}[H]
\centering
\small
\begin{tabular}{@{}lll@{}}
\toprule
 & \textbf{Datastore} & \textbf{Decentralized SQL} \\
\midrule
Shape      & columnar OLAP        & row OLTP (SQLite) \\
Scope      & aggregate / fleet    & per-user / project \\
Compute    & server-side (TEE)    & client + FHE on server \\
Privacy    & attested enclave     & E2E, keys client-held \\
Scale      & horizontal compute   & per-file, CRDT merge \\
Consensus  & Quasar (manifests)   & Quasar (keys/versions) \\
\bottomrule
\end{tabular}
\caption{Two engines, one platform: the OLAP/OLTP split re-cast along the trust
axis. Both ride S-Chain storage, K-Chain keys, and ZAP transport.}
\label{tab:twoengine}
\end{table}

\paragraph{The shared stack.} \textbf{S-Chain} (the Hanzo S3 / SeaweedFS fork)
is the decentralized storage tier for \emph{both} engines --- Datastore parts and
encrypted SQLite files are both content-addressed objects. \textbf{K-Chain}
custodies the post-quantum keys. \textbf{Z-Chain} provides FHE/ZK compute.
\textbf{Quasar} provides post-quantum, leaderless finality for the commitments of
both. And \textbf{ZAP} --- the Cap'n-Proto-derived transport that already
replaced gRPC --- is the single wire for client-to-engine queries, engine-to-S3
part transfer, and engine-to-consensus coordination, so the entire platform is
one transport deep rather than the usual stack of HTTP, gRPC, and a ZooKeeper
wire. The endpoint of the program in this paper is precisely that: no internal
HTTP, no gRPC, no ZooKeeper --- ZAP and Quasar, end to end, over decentralized
S-Chain storage.
